\subsection{Current Computational Approaches and Limitations}
\label{sec2.2: current approaches}
% State of the Art in Uncertainty Quantification for Biological Prediction
\noindent The problem of predicting which phage infects which host bacteria has been studied computationally for over a decade. The host 
prediction methods can be divided into categories, these methods progress from simple genomic statistics to deep 
representation learning. Here we will summarize each family of methods and their underlying mathematical assumptions along with
the common limitations that motivate us for this work.


\subsubsection{Sequence Similarity Approaches}
\label{sec2.2.1: SSA}
\noindent The earliest host prediction tools relied on sequence alignment between a phage and bacterial genomes, assuming that similarity 
between the two indicates infection compatibility. This was motivated by the known transfer of genetic material between the 
phages and their hosts through horizontal gene transfer.

\vspace{0.15cm}
\noindent One of the tools, HostPhinder \cite{villarroel2016hostphinder}, works by building a reference database of phage genomes with the
known hosts. Then for a new phage in question, it assigns the host that is the host of its nearest neighbour under a k-mer similarity
metric. Mathematically, let $f(g) \in \mathbb{R}^{4^k}$ be the normalised k-mer frequency vector of a genome $g$, then its 
predicted host is:

$$\hat{y} = \operatorname*{argmax}_{y \in \mathcal{Y}} \left\{ \operatorname{sim}(f(g_{\text{new}}), f(g_y)) \right\}$$

\noindent where $g(y)$ is the reference phage genome with a knwon host $y$ and $\operatorname*{sim}$ is a similarity measure.

\vspace{0.15cm}
\noindent Another method, WIsH \cite{galiez2017wish}, uses a probabilistic approach by modelling the genome of each candidate
host as a varaible order Markov chain and computing the likelihood of the new phage genome under each host model. The predicted
host maximises this likelihood. 

\vspace{0.15cm}
\noindent Despite the fact that these methods are relitively simpler, they perform reasonably well when they have a phage
reference that is closely related. But for a new phage that is not very common and is only a distant relative of the phages
already there in the database, they degrade sharply. And this scenario is very common in metagenomic settings \cite{coclet2021global}.


\subsubsection{Protein Content Approaches}
\label{sec2.2.2: PCA}
\noindent Instead of comparing and aligning the whole genomes, like the methods in \ref{sec2.2.1: SSA} these methods 
break down the phage genome into its constituent proteins and then analyze their specific features from their functional
or sequence properties. This approach is motivated by the fact that the proteins responsible for infecting a host are much
more evolutionary conserved than the rest of the genome \cite{coutinho2021rafah}.

\vspace{0.15cm}
\noindent RaFAH \cite{coutinho2021rafah} (Random Forest with protein Amino acid compositions for phage Host prediction)
works by clustering all the known phage proteins into protein clusters using MMseqs2 \cite{steinegger2017mmseqs2}.
It computes the amino acid composition vector of each cluster's representative sequence amd then trains a random forest
classifier to predict the host genus from the protein cluster membership profile of the new phage. Let $x \in \{0,1\}^M$
encode which of the $M$ protein clusters are present in a given phage, thenn random forest $f: \{0,1\}^M \to \varDelta(\mathcal{Y})$ 
maps the binary feature vector to a probability dostribution over the hosts $\mathcal{Y}$. 

\vspace{0.15cm}
\noindent RaFAH is also relevant to this report, since we also work with similar protein based cluster representation
and the contamination problem discussed in \todo{add section reference} also arises exactly from this protein-cluster 
construction.

\vspace{0.15cm}
\noindent These techniques are more robust to sequence divergence than the sequence alignment methods \ref{sec2.2.1: SSA}, but they rely on the availibility of explicit functional
annotations, which are missing for a large percentage of the already discovered phage genomes.



\subsubsection{Protein Language Models and Embedding Based Approaches}
\label{sec2.2.3: PLMs}
\noindent The latest host prediction methods use protein language models (pLMs) \cite{elnaggar2022prottrans, flamholz2024large}.
Similar to how large language models (LLMs) are trained on massive text datasets, pLMs are neural networks trained on 
millions of protein sequences through self-supervised learning.

\vspace{0.15cm}
\noindent For a protein sequence $a = (a_1, \dots, a_n) \in \mathcal{A}^n$, where $\mathcal{A}$ represents the standard
amino acid alphabet of size 20, a pLM maps the sequence to a dense vector representation:
$$z = \Phi(a) \in \mathbb{R}^d,$$
where $d$ is the embedding dimension (typically 1024 or larger). This embedding vector, $z$, captures the sequence's
biochemical, structural, and evolutionary properties. A primary foundational model in this domain is ProtTrans/ProtT5 
\cite{elnaggar2022prottrans}, which is trained on hundreds of millions of UniRef protein sequences.

\vspace{0.15cm}
\noindent Several tools for phage protein protein annotation and host prediction have been developed using these embeddings.
For instance, EMPATHI \cite{boulay2025empathi} uses a heirarchical classification pipeline to assigns functioal labels 
to phage proteins. It embeds the sequences using ProtT5 and performs nearest neighbour retreival witihin the embedding space,
mapping them to a biological function ontology. SUBLYME \cite{boulay2026phalp} then extends this methodology for metagenomics
which enables lysin discovery directly from environmemtal samples. These tools are directly involved in this work as rhe 
score matrix that is used as an input to this conformal framework is generated by a pipeline which is built on these annotations.

\vspace{0.15cm}
\noindent For host prediction, we can see in \cite{gonzales2023protein} that protein embeddings actually improve the 
predictions for the phage host interactions compared to the sequence-similarity baselines \ref{sec2.2.1: SSA}. In 
\cite{du2026high} Du et al. show that through some key protein embeddings LLMs can acheive high resolution host assignments.
Such results support the theory that embedding represenations contain information that is sufficeint for an accurate 
host-prediction. But at the same time, all these methods output point predictions, a single score ot a single host label,
and they do not provide us with any measure of uncertainity.

\subsubsection{Limitations of Existing Methods}






\vspace{4cm}


 \noindent - Brief overview of existing approaches in machine learning\\
 % Bayesian methods, conformal prediction, calibration methods
- Why Bayesian methods are unsuitable here: require distributional assumptions, computationally expensive for large-scale genomic data, posteriors depend heavily on prior choice\\
- Why calibration methods 
% (Platt scaling, temperature scaling) 
are insufficient: they produce calibrated probabilities but no formal finite-sample guarantee\\
- Why conformal prediction is the right choice: distribution-free, finite-sample guarantee, no assumptions beyond exchangeability, computationally lightweight




% \subsection{The Critical Deficiency: Failure of Uncertainty Quantification}
% \label{sec2.3:failure_of_uq}
% \noindent Despite their empirical success, all three generations of computational tools share a fundamental, systemic flaw: 
%they operate exclusively as \textbf{point predictors}. For an input phage feature vector $\mathbf{x} \in \mathcal{X}$, these 
%algorithms output a soft score vector $\mathbf{s}(\mathbf{x}) = \left(s_1(\mathbf{x}), \dots, s_K(\mathbf{x})\right) \in [0, 1]^K$, 
%where $s_k(\mathbf{x})$ is interpreted as a heuristic confidence or empirical probability that the phage infects host $k$. 
%Crucially, these scores completely lack rigorous **Uncertainty Quantification (UQ)**. They fail to decouple 
%\textit{aleatoric uncertainty} (inherent stochasticity in the biological infection process) from \textit{epistemic uncertainty} 
%(lack of model knowledge due to sparse training data) \cite{ochoarivera2024conformal}. 

% In clinical phage therapy, executing choices based on raw point predictions introduces severe risks. If a model yields 
%$s_1(\mathbf{x}) = 0.51$ for Host 1 and $s_2(\mathbf{x}) = 0.49$ for Host 2, a point-prediction rule forces the selection of 
%Host 1, despite the model being highly uncertain. To mitigate this risk, we analyze why standard statistical machine learning
%approaches fail to provide satisfactory guarantees.

% \subsubsection{Why Bayesian Formulations Are Unsuitable}
% \noindent A classical mathematical approach to UQ is the Bayesian paradigm, which targets the posterior predictive distribution 
%$P(Y \mid X, \mathcal{D})$. While theoretically elegant, Bayesian methods are poorly suited for large-scale phage-host modeling 
%for three reasons:
% \begin{enumerate}
%     \item \textbf{Computational Intractability:} Evaluating the posterior distribution over high-dimensional protein embedding 
            % spaces requires complex Markov Chain Monte Carlo (MCMC) sampling or Variational Inference (VI). This is computationally
            % prohibitive when scaled to massive matrices containing millions of bacterial and viral proteins \cite{steinegger2017mmseqs2}.
%     \item \textbf{Prior Sensitivity:} In high-dimensional, sparse biological matrices, the calculated posterior distributions depend 
            % heavily on arbitrary choices of prior distributions, which lack biological or statistical justification.
%     \item \textbf{Misspecification Vulnerability:} Bayesian guarantees are strictly asymptotic and rely on the structural assumption
            % that the true data-generating distribution lies within the chosen model class ($\mathcal{M}$-closed setting). If the 
            % underlying model architecture is misspecified, the resulting credible intervals lose their validity.
% \end{enumerate}

% \subsubsection{Why Standard Calibration Methods Are Insufficient}
% \noindent Alternative frequentist methods focus on post-hoc calibration techniques, such as Platt Scaling or Temperature Scaling. 
% These methods solve an optimization problem over an independent validation set to ensure that the output scores match empirical 
% frequencies, meaning:
% \begin{equation}
%     P\left(Y = k \mid s_k(X) = \hat{p}\right) \longrightarrow \hat{p} \quad \text{as } n \to \infty
% \end{equation}
% While calibration ensures that a confidence score of $0.90$ aligns with a $90\%$ long-run empirical success rate on average, 
% it suffers from severe limitations:
% \begin{enumerate}
%     \item \textbf{No Finite-Sample Guarantees:} Calibration metrics (such as the Expected Calibration Error, ECE) only provide 
            % asymptotic regularities. They offer no statistical safeguards when working with small, finite datasets or highly 
            % specialized clinical isolates.
%     \item \textbf{Retention of the Point Prediction Paradigm:} Calibration alters the scale of the scalar outputs but does not 
            % change the core prediction format. It does not provide a mathematical mechanism to construct a discrete predictive 
            % set $C(\mathbf{x}) \subseteq \mathcal{Y}$ that controls the exact instance-level error rate.
% \end{enumerate}

% This methodological gap demonstrates the need for a framework that is computationally lightweight, completely distribution-free, 
% and capable of rendering exact, finite-sample statistical guarantees. This directly motivates our development of a **conformal 
% prediction** pipeline.
